\documentclass[12pt,a4paper]{article}
\usepackage[utf8]{inputenc}
\usepackage[T1]{fontenc}
\usepackage[spanish]{babel}
\usepackage{amsmath, amssymb}
\usepackage{geometry}
\usepackage[explicit]{titlesec}
\usepackage{enumitem}
\usepackage{xcolor}
\usepackage{mdframed}
\usepackage{booktabs}
\usepackage{fancyhdr}
\usepackage{hyperref}

\geometry{margin=2.2cm, top=2.5cm, bottom=2.5cm, headheight=15pt}

% Colors
\definecolor{mainblue}{RGB}{25, 75, 140}
\definecolor{lightblue}{RGB}{220, 235, 252}
\definecolor{accentgray}{RGB}{245,246,248}
\definecolor{rulegray}{RGB}{180,190,210}
\definecolor{warnorange}{RGB}{255,243,224}
\definecolor{warnorangeline}{RGB}{210,120,0}

% Title formatting
\titleformat{\section}[block]
  {\Large\bfseries\color{mainblue}}
  {}
  {0em}
  {\colorbox{lightblue}{\parbox{\dimexpr\linewidth-2\fboxsep}{\thesection\quad #1}}}
  [\vspace{2pt}\color{rulegray}\hrule height 0.5pt\vspace{4pt}]

\titleformat{\subsection}[hang]
  {\normalsize\bfseries\color{mainblue}}
  {\thesubsection.}
  {0.5em}
  {}

\titlespacing*{\section}{0pt}{14pt}{6pt}
\titlespacing*{\subsection}{0pt}{10pt}{4pt}

% Box for key formulas
\newmdenv[
  backgroundcolor=accentgray,
  linecolor=mainblue,
  linewidth=1pt,
  topline=true, bottomline=true, leftline=true, rightline=false,
  innerleftmargin=10pt, innerrightmargin=10pt,
  innertopmargin=6pt, innerbottommargin=6pt,
  leftmargin=0pt, rightmargin=0pt,
  skipabove=6pt, skipbelow=6pt
]{formulabox}

% Box for warnings / important notes
\newmdenv[
  backgroundcolor=warnorange,
  linecolor=warnorangeline,
  linewidth=1.2pt,
  topline=true, bottomline=true, leftline=true, rightline=false,
  innerleftmargin=10pt, innerrightmargin=10pt,
  innertopmargin=6pt, innerbottommargin=6pt,
  leftmargin=0pt, rightmargin=0pt,
  skipabove=6pt, skipbelow=6pt
]{warnbox}

% Header/footer
\pagestyle{fancy}
\fancyhf{}
\fancyhead[L]{\small\color{mainblue}\textbf{Apuntes de Circuitos Eléctricos}}
\fancyhead[R]{\small\color{mainblue}\thepage}
\renewcommand{\headrulewidth}{0.4pt}
\renewcommand{\headrule}{\color{rulegray}\hrule width\headwidth height\headrulewidth}

\hypersetup{colorlinks=true, linkcolor=mainblue, urlcolor=mainblue}

\begin{document}

% Title
\begin{center}
  {\huge\bfseries\color{mainblue} Apuntes de Circuitos Eléctricos}\\[0.4em]
  {\large\color{gray} Temas 1, 3 y 4}\\[0.3em]
  \color{rulegray}\rule{0.6\linewidth}{1pt}
\end{center}
\vspace{1em}

% ─────────────────────────────────────────────────────────────────────────────
\section*{Glosario de Símbolos y Conceptos}

\textbf{\color{mainblue}Símbolos y variables:}

\begin{itemize}[leftmargin=2em, itemsep=3pt, parsep=0pt]
  \item[$\tau$] \textbf{Constante de tiempo:} caracteriza la velocidad del transitorio. En primer orden, $\tau = R_{eq}C$ (RC) o $\tau = L/R_{eq}$ (RL).
  \item[$\omega_n$] \textbf{Pulsación natural:} velocidad característica del sistema sin amortiguamiento. No confundir con $\omega_d$.
  \item[$\omega_d$] \textbf{Pulsación amortiguada:} frecuencia real a la que oscila el sistema. $\omega_d = \omega_n\sqrt{1-\xi^2}$.
  \item[$\omega_0$] \textbf{Pulsación de resonancia:} frecuencia donde reactancias se anulan. $\omega_0 = 1/\sqrt{LC}$.
  \item[$\omega_c$] \textbf{Pulsación de corte (-3\,dB):} límite entre banda pasante y atenuada. Solo fórmula directa en orden 1.
  \item[$\xi$] \textbf{Factor de amortiguamiento:} controla oscilación ($\xi < 1 \Rightarrow$ oscila, $\xi \geq 1 \Rightarrow$ no). $\xi = 1/2Q$.
  \item[$Q$] \textbf{Factor de calidad:} relación energía almacenada / disipada. $Q = \omega_0/B$. Alto $Q \Rightarrow$ pico resonante estrecho.
  \item[$B$] \textbf{Ancho de banda:} rango de frecuencias a $-3$\,dB donde cae menos de 3\,dB. En rad/s.
  \item[$R_{eq}$] \textbf{Resistencia equivalente:} se obtiene apagando fuentes independientes. Determina $\tau$ y amortiguamiento.
\end{itemize}

\medskip
\textbf{\color{mainblue}Conceptos principales:}

\begin{itemize}[leftmargin=1.5em, itemsep=8pt, parsep=0pt]
  \item \textbf{Polinomio Característico:} el denominador de la función de transferencia $H(s)$. Sus raíces (polos) determinan si el sistema es estable y qué tipo de respuesta tiene (amortiguada, oscilatoria, etc.).
  
  \item \textbf{Régimen Transitorio:} la respuesta inicial tras un cambio. Decae exponencialmente (o oscilatoriamente según $\xi$) con constante de tiempo $\tau$ o $\omega_n$. Es la parte $y_h(t)$ que desaparece con el tiempo.
  
  \item \textbf{Régimen Permanente:} la respuesta en $t \to \infty$, una vez que desaparece el transitorio. El sistema alcanza el nuevo equilibrio. Es la solución particular $y_p(t)$.
  
  \item \textbf{Resonancia (RLC Serie):} cuando la reactancia neta es cero ($\Im(Z) = 0$). Ocurre a $\omega_0 = 1/\sqrt{LC}$. La impedancia es mínima (solo $R$), corriente máxima, tensiones en $L$ y $C$ amplificadas por factor $Q_s$.
  
  \item \textbf{Impedancia de Thévenin} ($Z_{Th}$): la resistencia equivalente vista desde la carga, con fuentes independientes apagadas. Determina el acoplamiento y adaptación de impedancias ($Z_L = Z_{Th}^*$ para máxima potencia).
\end{itemize}

% ─────────────────────────────────────────────────────────────────────────────
\section{Régimen Transitorio en Circuitos}

\textbf{Objetivo:} Analizar la respuesta de los circuitos en cualquier instante de tiempo,
especialmente durante el período transitorio (desde el cambio hasta el régimen permanente).

% ── 1.1 ──────────────────────────────────────────────────────────────────────
\subsection{Conceptos Clave}

\begin{itemize}[leftmargin=1.5em, itemsep=4pt]

  \item \textbf{Respuesta completa:}
  \[
    y(t) = y_h(t) + y_p(t)
  \]
  donde $y_h(t)$ es la \textbf{respuesta homogénea} (transitorio: parte que decae con $\tau$) y $y_p(t)$ la
  \textbf{respuesta particular} (régimen permanente: solución final estable).

  \item \textbf{Condiciones iniciales:} $v_C$ e $i_L$ no pueden variar bruscamente:
  \[
    v_C(0^+) = v_C(0^-), \qquad i_L(0^+) = i_L(0^-)
  \]

  \item \textbf{Transformada de Laplace (unilateral):}
  \[
    F(s) = \mathcal{L}[f(t)] = \int_0^\infty f(t)\,e^{-st}\,dt
  \]
  Propiedades útiles:
  \begin{align*}
    \mathcal{L}[f'(t)] &= sF(s) - f(0^-)\\
    \mathcal{L}\!\left[\int_{-\infty}^t f(\tau)\,d\tau\right] &= \frac{F(s)}{s}
      + \frac{f^{(-1)}(0^-)}{s}
  \end{align*}
\end{itemize}
\clearpage %

% ── 1.2 ──────────────────────────────────────────────────────────────────────
\subsection{Circuito de Primer Orden — Algoritmo Universal}

En cualquier circuito de primer orden (un solo elemento reactivo) la variable $x(t)$
(tensión o corriente) sigue siempre:

\begin{formulabox}
\[
  \boxed{x(t) = X_\infty + \bigl(X_{0^+} - X_\infty\bigr)\,e^{-t/\tau}}
\]
\end{formulabox}

Los tres ingredientes se calculan \textbf{en este orden}:

\begin{enumerate}[leftmargin=2em, itemsep=6pt]

  \item \textbf{Condición inicial} $X_{0^+}$: resolver el circuito en $t = 0^-$
  (régimen CC previo al cambio) usando las equivalencias de CC:
  \begin{formulabox}
    \textbf{En CC} ($t = 0^-$ o $t \to \infty$):\\[4pt]
    \quad Condensador $\longrightarrow$ \textbf{circuito abierto} \quad
    (no pasa corriente, acumula tensión)\\
    \quad Bobina $\longrightarrow$ \textbf{cortocircuito} \quad
    (no cae tensión, acumula corriente)\\[4pt]
    La tensión/corriente obtenida en $t = 0^-$ es el valor inicial del elemento reactivo
    en $t = 0^+$ (la condición no cambia bruscamente).
  \end{formulabox}

  \item \textbf{Valor final} $X_\infty$: resolver el circuito en $t \to \infty$
  (nuevo régimen CC tras el cambio) con las mismas equivalencias del punto anterior.

  \item \textbf{Constante de tiempo} $\tau$ (velocidad del transitorio): apagar todas las \emph{fuentes independientes}
  (tensión $\to$ cortocircuito, corriente $\to$ circuito abierto) y calcular la resistencia
  equivalente $R_{eq}$ (resistencia vista desde el elemento reactivo):
  \[
    \tau = R_{eq}\,C \quad \text{(RC)} \qquad \tau = \frac{L}{R_{eq}} \quad \text{(RL)}
  \]
\end{enumerate}
\clearpage %
% ── 1.3 ──────────────────────────────────────────────────────────────────────
\subsection{Circuito de Segundo Orden — Ecuación Característica}

La ecuación diferencial de un circuito de segundo orden tiene siempre la forma:
\[
  s^2 + 2\xi\omega_n\,s + \omega_n^2 = 0
\]

\begin{warnbox}
\textbf{\color{warnorangeline}Atención:} el factor de amortiguamiento $\xi$ (determina si hay oscilación: $\xi < 1 \Rightarrow$ oscila, $\xi \geq 1 \Rightarrow$ no) y la pulsación natural $\omega_n$ (velocidad característica del sistema) \textbf{no} se leen
directamente del circuito con una fórmula fija. El procedimiento correcto es:
\begin{enumerate}[leftmargin=2.5em, itemsep=3pt]
  \item Obtener la función de transferencia $H(s)$ o la impedancia en el dominio de
        Laplace (con condiciones iniciales nulas).
  \item Identificar el polinomio característico en el denominador.
  \item Igualar coeficientes con la forma estándar para extraer $\xi$ y $\omega_n$.
\end{enumerate}
La expresión $\xi = \tfrac{R}{2}\sqrt{C/L}$ \textbf{solo vale para el RLC serie básico}.
En circuitos mixtos siempre hay que hacer el paso 1--3.
\end{warnbox}

\medskip
\textbf{Tipos de respuesta según $\xi$:}

\begin{center}
\renewcommand{\arraystretch}{1.4}
\begin{tabular}{llll}
\toprule
\textbf{Condición} & \textbf{Tipo} & \textbf{Raíces} & \textbf{Comportamiento}\\
\midrule
$0 < \xi < 1$ & Subamortiguado    & Complejas conjugadas & Oscilatoria amortiguada\\
$\xi = 1$     & Amortiguamiento crítico & Reales e iguales & Más rápida sin oscilación\\
$\xi > 1$     & Sobreamortiguado  & Reales y distintas   & Lenta sin oscilación\\
$\xi = 0$     & No amortiguado    & Imaginarias puras    & Oscilación permanente\\
\bottomrule
\end{tabular}
\end{center}

Para el caso subamortiguado ($0 < \xi < 1$):
\[
  s_{1,2} = -\omega_n\xi \pm j\,\omega_d, \qquad \omega_d = \omega_n\sqrt{1-\xi^2}
\]

\begin{warnbox}
\textbf{\color{warnorangeline}Atención a la oscilación en tiempo:}\\[4pt]
En el dominio del tiempo, el seno/coseno oscila a la pulsación amortiguada $\omega_d$ (frecuencia real de oscilación cuando hay amortiguamiento),
\textbf{NO} a la $\omega_n$ (pulsación natural: la que habría sin amortiguamiento). En exámenes suelen pedir
dibujar la señal o escribirla, y el fallo más típico es poner $\omega_n$ dentro
del seno.
\end{warnbox}

\vspace{4pt}
\textbf{Nota geométrica:} En el plano $s$, la distancia del polo al origen es
$|s_{1,2}| = \omega_n$ (pulsación natural: magnitud del polo), no $\omega_d$ (pulsación amortiguada: solo la parte imaginaria).
\clearpage %
% ── 1.4 ──────────────────────────────────────────────────────────────────────
\subsection{Equivalentes en el Dominio de Laplace (condiciones iniciales no nulas)}

Los modelos siguientes corresponden al \textbf{equivalente serie}: la impedancia del
elemento queda en serie con una fuente de tensión que representa la condición inicial.
Se aplica directamente la LKT en el dominio $s$.

\begin{formulabox}
\textbf{Resistencia:}
\[
  V_R(s) = R\cdot I_R(s)
\]
\textbf{Condensador} — impedancia $1/sC$ en serie con generador de tensión $v_C(0^-)/s$:
\[
  V_C(s) = \frac{1}{sC}\,I_C(s) + \frac{v_C(0^-)}{s}
\]
\textbf{Bobina} — impedancia $sL$ en serie con generador de tensión $L\,i_L(0^-)$:
\[
  V_L(s) = sL\cdot I_L(s) - L\cdot i_L(0^-)
\]
El signo negativo en la bobina sigue la convención: si $i_L(0^-)>0$, la fuente
serie actúa a favor de la corriente, reduciendo la caída neta de tensión.
\end{formulabox}

\begin{warnbox}
\textbf{\color{warnorangeline}Importante — modelo serie, no paralelo:}\\[4pt]
Existe un modelo dual equivalente (fuente de \emph{corriente} en \emph{paralelo} con
la admitancia), pero \textbf{estos apuntes usan siempre el modelo serie}. Al plantear
la malla en el dominio $s$ con LKT, las condiciones iniciales aparecen como caídas
de tensión adicionales (fuentes en serie); no es necesario convertir fuentes ni
trabajar con admitancias.
\end{warnbox}
\clearpage %
% ─────────────────────────────────────────────────────────────────────────────
\section{Transferencia de Potencia y Adaptación de Impedancias}

\textbf{Objetivo:} Determinar las condiciones para máxima transferencia de potencia
de un generador a una carga, y estudiar el efecto de redes intermedias.

% ── 2.1 ──────────────────────────────────────────────────────────────────────
\subsection{Conceptos Clave}

\begin{itemize}[leftmargin=1.5em, itemsep=4pt]
  \item \textbf{Adaptación de impedancias:} condición en la que la máxima potencia se transfiere (ocurre cuando $Z_L = Z_{Th}^*$, el conjugado de la impedancia de Thévenin).

  \item \textbf{Teorema de Everitt:} si entre generador y carga se intercala una red
        \emph{LC pura} (sin pérdidas) y existe adaptación en un punto, existe adaptación
        en todos los puntos de la red y la potencia máxima se transfiere íntegramente
        hasta la carga. Corolario útil en examen: si hay adaptación a la entrada de
        la red LC, la hay también a la salida, y por tanto $\alpha_T = 0\,\text{dB}$.
        Una red \emph{RLC} (disipativa) \textbf{rompe} esta condición.

  \item \textbf{Pérdidas de inserción ($\alpha_I$):} medida logarítmica de la potencia
        que se pierde al insertar una red entre generador y carga.
\end{itemize}

% ── 2.2 ──────────────────────────────────────────────────────────────────────
\subsection{Cálculo de $Z_{Th}$ — Paso Previo Imprescindible}

\begin{warnbox}
\textbf{\color{warnorangeline}Regla de oro:} para calcular la impedancia de Thévenin $Z_{Th}$ (resistencia equivalente vista desde la carga), \textbf{apagar todas las fuentes independientes}
(fuentes de tensión $\to$ cortocircuito; fuentes de corriente $\to$ circuito abierto)
y calcular la impedancia entre los bornes de la carga.\\[4pt]
\textbf{Con fuentes dependientes:} no se puede simplificar directamente. Hay que aplicar
el \emph{método de la fuente de prueba}: conectar $V_{test}$ (o $I_{test}$) en los
bornes con las fuentes independientes apagadas y calcular
$Z_{Th} = V_{test}/I_{test}$.
\end{warnbox}
\clearpage %
% ── 2.3 ──────────────────────────────────────────────────────────────────────
\subsection{Máxima Transferencia de Potencia}

Dado un equivalente Thévenin $\bigl(E_{Th},\;Z_{Th} = a + jb\bigr)$ (impedancia interna del generador) y una carga
$Z_L = x + jy$:

\begin{formulabox}
\textbf{Condición general (carga compleja):}
\[
  Z_L = Z_{Th}^* = a - jb
\]
\textbf{Potencia máxima:}
\[
  P_{L,\max} = \frac{|E_{Th}|^2}{8\,\Re\{Z_{Th}\}} = \frac{|E_{Th}|^2}{8a}
\]

\begin{warnbox}
\textbf{\color{warnorangeline}$\star\star\star$ CRUCIAL $\star\star\star$}\\[4pt]
El denominador \textbf{8} en la fórmula anterior \textbf{SOLO} se usa si $|E_{Th}|$
es la \emph{Amplitud Máxima} (pico) del generador (p.ej., el 4 en $4\cos(\omega t)$).
\emph{Si en el examen te dan el dato con un voltímetro/amperímetro, te están dando
el Valor Eficaz (RMS)}. En ese caso, la fórmula lleva un \textbf{4}:
\[
  P_{L,\max} = \frac{|E_{Th,\mathrm{ef}}|^2}{4\,\Re\{Z_{Th}\}}
\]
Muchísima gente suspende por usar el 8 con valores eficaces y perder la mitad
de la potencia.
\end{warnbox}

\vspace{6pt}
\textbf{Caso particular (carga resistiva pura, $Z_L = R_L$):}
\[
  R_L = |Z_{Th}| = \sqrt{a^2 + b^2}
  \qquad\text{(potencia máxima menor que el caso general)}
\]
\end{formulabox}
\clearpage %
% ── 2.4 ──────────────────────────────────────────────────────────────────────
\subsection{Unidades de Transmisión y Pérdidas de Inserción}

\begin{itemize}[leftmargin=1.5em, itemsep=6pt]

  \item \textbf{Decibelio (dB):} escala logarítmica para comparar magnitudes (potencia, tensión, corriente). Fórmula:
  \[
    X\,[\text{dB}] = 10\log_{10}\!\frac{P_1}{P_2}
    = 20\log_{10}\!\frac{V_1}{V_2} = 20\log_{10}\!\frac{I_1}{I_2}
  \]
  \begin{warnbox}
  \textbf{\color{warnorangeline}Condición de validez:} las formas con tensiones
  ($20\log V_1/V_2$) y con corrientes ($20\log I_1/I_2$) \textbf{solo son válidas si
  ambas magnitudes se miden sobre la misma impedancia}. En caso contrario, la
  expresión correcta es siempre la de potencias: $10\log(P_1/P_2)$.
  \end{warnbox}

  \item \textbf{Neper (Np):}
  \[
    Y\,[\text{Np}] = \frac{1}{2}\ln\frac{P_1}{P_2},
    \qquad 1\,\text{Np} = 20\log_{10}(e) \approx 8{,}69\,\text{dB}
  \]

  \item \textbf{Pérdidas de inserción:}
  \[
    \alpha_I\,[\text{dB}]
    = 10\log\frac{P_{20}}{P_2}
    = 20\log\frac{|V_{20}|}{|V_2|}
    = 20\log\frac{|I_{20}|}{|I_2|}
  \]

  \begin{warnbox}
  \textbf{\color{warnorangeline}Definición exacta de $P_{20}$ y $P_2$}
  (confusión frecuente en examen):\\[4pt]
  \quad $P_{20}$: potencia entregada a la carga cuando está conectada
  \textbf{directamente} al generador, \emph{sin} ninguna red intermedia.\\
  \quad $P_2$: potencia entregada a la misma carga \textbf{con} la red intermedia
  conectada.\\[4pt]
  Si la red es LC pura, $P_{entrada} = P_{salida}$, luego las Pérdidas de
  Transmisión son $\alpha_T = 0\,\text{dB}$. Las pérdidas de inserción ($\alpha_I$)
  casi nunca son 0, de hecho suelen ser negativas (ganancia) porque la red se
  pone precisamente para adaptar lo que antes estaba desadaptado.
  \end{warnbox}

\end{itemize}
\clearpage %
% ─────────────────────────────────────────────────────────────────────────────
\section{Respuesta en Frecuencia. Filtrado}

\textbf{Objetivo:} Analizar cómo varía el comportamiento de un circuito (amplitud
y fase) con la frecuencia; aplicarlo al diseño e identificación de filtros.

% ── 3.1 ──────────────────────────────────────────────────────────────────────
\subsection{Función de Transferencia y Respuesta en Frecuencia}

\begin{itemize}[leftmargin=1.5em, itemsep=4pt]
  \item $H(s) = Y(s)/X(s)$ con condiciones iniciales nulas.
  \item \textbf{Polos:} raíces del denominador (donde $H(s)\to\infty$). Determinan estabilidad y forma de respuesta. Se representan con \textbf{cruces (x)} en el plano $s$.
  \item \textbf{Ceros:} raíces del numerador (donde $H(s) = 0$). Indican atenuaciones a frecuencias específicas. Se representan con \textbf{círculos ($\circ$)} en el plano $s$.
  \item \textbf{Respuesta en frecuencia:} cómo cambian ganancia y fase con la frecuencia. Se calcula sustituyendo $s = j\omega$ en $H(s)$: el módulo $|H(j\omega)|$ da la ganancia, la fase $\angle H(j\omega)$ el desfase.
  \item \textbf{Resonancia:} las reactancias de bobina y condensador se anulan, impedancia mínima, corriente máxima.
  \[
    \omega_0 = \frac{1}{\sqrt{LC}} \quad\text{(pulsación de resonancia: donde se anulan reactancias)}, \quad Q = \frac{\omega_0}{B}
  \]
  donde $Q$ = factor de calidad (relación entre energía almacenada y disipada, alto $Q$ $\Rightarrow$ resonancia estrecha)
  y $B$ = ancho de banda a $-3$\,dB en rad/s (rango de frecuencias donde cae menos de 3\,dB).
\end{itemize}

% ── 3.2 ──────────────────────────────────────────────────────────────────────
\subsection{Identificación de Filtros por Análisis Asintótico}

Antes de calcular $H(s)$, se puede razonar el tipo de filtro evaluando los casos
límite directamente sobre el \emph{circuito} con estas equivalencias:

\begin{formulabox}
\renewcommand{\arraystretch}{1.5}
\begin{tabular}{lll}
\toprule
\textbf{Elemento} & $\omega \to 0$ \;(CC, baja freq.) & $\omega \to \infty$ (alta freq.)\\
\midrule
Condensador & Circuito \textbf{abierto} (bloquea DC) & Cortocircuito (deja pasar AC)\\
Bobina      & Cortocircuito (no resiste DC) & Circuito \textbf{abierto} (bloquea AC)\\
\bottomrule
\end{tabular}\\[8pt]
\textbf{Razonamiento:} Si la señal \textbf{llega} a la salida en $\omega\to 0$ pero \textbf{no} en
$\omega\to\infty$ $\Rightarrow$ \textbf{Paso Bajo}.\\
Si \textbf{no llega} en $\omega\to 0$ pero \textbf{sí} en
$\omega\to\infty$ $\Rightarrow$ \textbf{Paso Alto}.\\
Si \textbf{no llega} en ninguno de los dos extremos $\Rightarrow$ \textbf{Paso Banda}.\\
Si \textbf{llega} en ambos extremos $\Rightarrow$ \textbf{Banda Eliminada}.
\end{formulabox}
\clearpage %
% ── 3.3 ──────────────────────────────────────────────────────────────────────
\subsection{Diagrama de Polos y Ceros — Reglas de Identificación}

\begin{formulabox}
\textbf{Ceros en el origen} ($s = 0$): $H(0) = 0$ $\Rightarrow$ la señal \textbf{no
pasa} a CC (baja freq.) $\Rightarrow$ Filtro \textbf{PA (Paso Alto) o PBo (Paso Banda)}.\\[4pt]
\textbf{Sin ceros finitos} (o ceros en $s\to\infty$): $H(0) \neq 0$ $\Rightarrow$
la señal \textbf{pasa} en CC $\Rightarrow$ Filtro \textbf{PB (Paso Bajo)}.\\[4pt]
\textbf{Ceros en el eje imaginario} ($s = \pm j\omega_0$, no en el origen):
$H(j\omega_0) = 0$ $\Rightarrow$ atenuación total a esa frecuencia $\Rightarrow$
Filtro \textbf{BE (Banda Eliminada)}.\\[4pt]
\textbf{Pendiente asintótica} en alta frecuencia: cada polo en exceso sobre los ceros
contribuye $-20$\,dB/década (cambio de pendiente).
\end{formulabox}

\medskip
\textbf{Relación entre polos y parámetros del sistema (caso subamortiguado):}
Si el polo es $s = -\sigma \pm j\omega_d$ (con $\sigma > 0$), entonces
\[
  \omega_n = \sqrt{\sigma^2 + \omega_d^2} \quad\text{(distancia del polo al origen: pulsación natural)}
\]
Atención: $\omega_n$ es la \textbf{magnitud} del polo, $\omega_d$ es solo la \textbf{parte imaginaria}.

% ── 3.4 ──────────────────────────────────────────────────────────────────────
\subsection{Filtros de Primer Orden}

\begin{formulabox}
\[
  H_{PB}(s) = k\cdot\frac{\omega_c}{s+\omega_c}, \qquad
  H_{PA}(s) = k\cdot\frac{s}{s+\omega_c}
\]
donde $\omega_c$ = pulsación de corte a $-3$\,dB. Se define por $|H(\omega_c)| = |H_{\max}|/\sqrt{2}$.\\[4pt]
\textbf{Fórmulas SOLO para filtros de primer orden} (un solo elemento reactivo: RC o RL):
\[
  \text{Circuito RC: }\omega_c = \frac{1}{RC} \qquad \text{Circuito RL: }\omega_c = \frac{R}{L}
\]
Para circuitos de segundo orden (RLC), igualar el denominador a $s^2 + \dfrac{\omega_0}{Q}\,s + \omega_0^2$.
\end{formulabox}
\clearpage %
% ── 3.5 ──────────────────────────────────────────────────────────────────────
\subsection{Filtros de Segundo Orden}

Forma estándar del denominador (polinomio característico):
$\displaystyle s^2 + \frac{\omega_0}{Q}\,s + \omega_0^2$

donde $Q = \dfrac{1}{2\xi}$ (factor de calidad, relacionado con amortiguamiento),
$\omega_0 = \dfrac{1}{\sqrt{LC}}$ (pulsación de resonancia), y $B = \dfrac{\omega_0}{Q}$ (ancho de banda a $-3$\,dB).

\begin{formulabox}
\textbf{Paso Bajo (PB):}
\[
  H_{PB}(s) = k\cdot
  \frac{\omega_0^2}{s^2 + \dfrac{\omega_0}{Q}\,s + \omega_0^2}
\]
\textbf{Paso Alto (PA):}
\[
  H_{PA}(s) = k\cdot
  \frac{s^2}{s^2 + \dfrac{\omega_0}{Q}\,s + \omega_0^2}
\]
\textbf{Paso Banda (PBo):}
\[
  H_{PBo}(s) = k\cdot
  \frac{\dfrac{\omega_0}{Q}\,s}{s^2 + \dfrac{\omega_0}{Q}\,s + \omega_0^2},
  \qquad B = \frac{\omega_0}{Q}
\]
\textbf{Banda Eliminada (BE):}
\[
  H_{BE}(s) = k\cdot
  \frac{s^2 + \omega_0^2}{s^2 + \dfrac{\omega_0}{Q}\,s + \omega_0^2}
\]
\end{formulabox}
\clearpage %
% ── 3.6 ──────────────────────────────────────────────────────────────────────
\subsection{Resonancia Serie (y caso dual paralelo)}

\textbf{Condición de resonancia:} parte imaginaria de la impedancia serie igual a cero,
con $\omega_0 = 1/\!\sqrt{LC}$. El análisis base siempre parte de un circuito
\textbf{RLC en serie} con generador de tensión.

\begin{formulabox}
\textbf{Resonancia Serie} (circuito RLC en serie con generador de tensión $V_g$):
Cuando la reactancia neta es cero ($\Im(Z) = 0$), la impedancia es mínima e igual a $R_T$ (solo resistencia):
\[
  I_{\max} = \frac{V_g}{R_T}, \qquad
  Q_s = \frac{\omega_0 L}{R_T} = \frac{1}{\omega_0 C R_T} \quad\text{(factor de calidad serie)}, \qquad
  |V_L| = |V_C| = Q_s\cdot|V_g|
\]
El factor $Q_s$ amplifica las tensiones en los elementos reactivos. A resonancia: corriente es máxima, tensiones en $L$ y $C$ están amplificadas.

\medskip
\textbf{Caso dual — Resonancia Paralelo} (RLC en paralelo con generador de corriente $I_g$):\\[4pt]
Cuando la susceptancia neta es cero ($\Im(Y) = 0$), la admitancia es mínima ($Y = 1/R_T$), tensión máxima:
\[
  V_{\max} = I_g\cdot R_T, \qquad
  Q_p = \omega_0 C R_T = \frac{R_T}{\omega_0 L} \quad\text{(factor de calidad paralelo)}, \qquad
  |I_L| = |I_C| = Q_p\cdot|I_g|
\]
Las corrientes en las ramas se amplifican por $Q_p$.
\end{formulabox}

\begin{warnbox}
\textbf{\color{warnorangeline}Regla de trabajo:} ante cualquier circuito RLC,
plantear siempre la \textbf{malla en serie} (LKT con generador de tensión).
El circuito paralelo es el dual matemático; si aparece en un problema,
convertir el generador de corriente en paralelo a su equivalente Thévenin
(generador de tensión en serie con resistencia equivalente) antes de aplicar las fórmulas de $Q_s$.
\end{warnbox}

\end{document}
